\documentclass[12pt]{article}
\usepackage[UTF8]{ctex}
\usepackage{amsmath,amssymb}
\usepackage{graphicx}
\usepackage{booktabs}
\usepackage{multirow}
\usepackage{multicol}
\usepackage{geometry}
\geometry{a4paper,left=2.5cm,right=2.5cm,top=2.5cm,bottom=2.5cm}

\title{yidoer的热学期末复习资料}
\author{by 元宝}
\date{ }

\begin{document}

\maketitle

%\tableofcontents

\section*{前言}
本资料基于信陽師範大学翟春洋老师的《热学》课程讲义（2025年10月-12月）整理，涵盖热学全部核心内容，包括基本概念、重要公式、图表及复习要点。资料按章节组织，力求全面系统，便于期末复习。

\section{第一章：热学导论}
\subsection{热学基本概念与研究方法}

热现象与热运动
\begin{itemize}
    \item 热现象：与物体冷热程度相关的物理性质变化
    \item 热运动：大量分子无规则运动的宏观体现
\end{itemize}

研究方法
\begin{itemize}
    \item 宏观描述（热力学方法）：基于实验定律
    \item 微观描述（统计物理）：从分子运动出发
\end{itemize}

热力学系统分类
\begin{itemize}
    \item 孤立系统：无能量、物质交换
    \item 封闭系统：有能量、无物质交换  
    \item 开放系统：有能量、物质交换
\end{itemize}

\subsection{平衡态与非平衡态}
平衡态条件
\begin{itemize}
    \item 热平衡：系统内部温度相等
    \item 力学平衡：无宏观粒子流
    \item 化学平衡：无相变发生
\end{itemize}
判断标准
是否存在热流与粒子流

\subsection{温度与热力学第零定律}
温度定义
分子热运动剧烈程度的量度

热力学第零定律
若A与B均与C热平衡，则A与B彼此热平衡

温标建立
\begin{align*}
    T &= \lim_{p_{tr} \to 0} \left(273.16 \cdot \frac{p}{p_{tr}}\right) \, \text{K} \\
    t_{90} &= \left(T_{90} - 273.15\right) \, ^\circ\text{C} \\
    t_F &= \left(\frac{9}{5} \frac{t}{^\circ\text{C}} + 32\right) \, ^\circ\text{F}
\end{align*}

\subsection{物态方程}
基本方程
\begin{align*}
    f(T, p, V) &= 0 \\
    pV &= vRT \quad (\text{理想气体})
\end{align*}

系数定义
\begin{align*}
    \kappa_T &= -\frac{1}{V} \left(\frac{\partial V}{\partial p}\right)_T \quad \text{(等温压缩系数)} \\
    \beta_p &= \frac{1}{V} \left(\frac{\partial V}{\partial T}\right)_p \quad \text{(体膨胀系数)} \\
    \alpha_V &= \frac{1}{p} \left(\frac{\partial p}{\partial T}\right)_V \quad \text{(压强系数)} \\
    p\alpha_V\kappa_T &= \beta_p
\end{align*}

混合理想气体
\begin{align*}
    pV &= (v_1 + v_2 + \cdots + v_n)RT \\
    p &= p_1 + p_2 + \cdots + p_n \quad \text{(道尔顿分压定律)}
\end{align*}

\subsection{物质的微观模型}
阿伏伽德罗常量
\[
N_A = 6.02 \times 10^{23} \, \text{mol}^{-1}
\]

分子力与分子热运动
分子吸引力与排斥力平衡距离为$r_0$

\subsection{理想气体微观描述的初级理论}
微观模型假设
\begin{itemize}
    \item 分子线度远小于间距
    \item 除碰撞外分子间无作用力  
    \item 碰撞为完全弹性碰撞
\end{itemize}

重要公式
\begin{align*}
    n_0 &= \frac{N_A}{22.4 \times 10^{-3}} = 2.7 \times 10^{25} \, \text{m}^{-3} \\
    \bar{\lambda} &= \frac{1}{\sqrt{2}n\sigma}, \quad \sigma = \pi d^2 \\
    \bar{Z} &= \sqrt{2}n\bar{v}\sigma \\
    \Gamma &= \frac{1}{6}n\bar{v} \\
    p &= \frac{1}{3}nm\overline{v^2} = \frac{2}{3}n\overline{\varepsilon_t} \\
    \overline{\varepsilon_t} &= \frac{3}{2}kT \\
    v_{\text{rms}} &= \sqrt{\frac{3kT}{m}} = \sqrt{\frac{3RT}{M_m}} \\
    p &= nkT
\end{align*}

\subsection{真实气体物态方程}
范德瓦尔斯方程
\[
\left(p + \frac{a}{V_m^2}\right)(V_m - b) = RT
\]

昂内斯方程
\[
pV = A + \frac{B}{V} + \frac{C}{V^2} + \cdots
\]

\section{第二章：分子动理学理论的平衡态理论}
\subsection{概率论基础}
基本定义
\begin{align*}
    P_i &= \lim_{N \to \infty} \frac{N_i}{N} \\
    \bar{u} &= \sum u_i P_i \\
    \overline{(\Delta u)^2} &= \overline{u^2} - (\bar{u})^2
\end{align*}

\subsection{麦克斯韦速率分布}
分布函数
\[
f(v)dv = 4\pi\left(\frac{m}{2\pi kT}\right)^{3/2} e^{-\frac{mv^2}{2kT}} v^2 dv
\]

三种特征速率
\begin{align*}
    v_p &= \sqrt{\frac{2kT}{m}} \quad \text{(最概然速率)} \\
    \bar{v} &= \sqrt{\frac{8kT}{\pi m}} \quad \text{(平均速率)} \\
    v_{\text{rms}} &= \sqrt{\frac{3kT}{m}} \quad \text{(方均根速率)}
\end{align*}

\subsection{麦克斯韦速度分布}
速度空间分布
\[
f(v_x, v_y, v_z)dv_x dv_y dv_z = \left(\frac{m}{2\pi kT}\right)^{3/2} e^{-\frac{m(v_x^2+v_y^2+v_z^2)}{2kT}} dv_x dv_y dv_z
\]

\subsection{玻尔兹曼分布}
等温大气公式
\begin{align*}
    p(z) &= p(0)e^{-\frac{M_m g z}{RT}} \\
    n(z) &= n(0)e^{-\frac{M_m g z}{RT}}
\end{align*}

玻尔兹曼分布
\[
\frac{N_1}{N_2} = e^{-\frac{\varepsilon_1 - \varepsilon_2}{kT}}
\]

\subsection{能量均分定理}
自由度
\begin{itemize}
    \item 单原子分子：3个平动自由度
    \item 刚性双原子分子：5个自由度（3平动+2转动）
    \item 刚性多原子分子：6个自由度（3平动+3转动）
\end{itemize}

能量均分公式
\begin{align*}
    \bar{\varepsilon} &= \frac{i}{2}kT \\
    U_m &= \frac{i}{2}RT \\
    C_{V,m} &= \frac{i}{2}R \\
    C_{p,m} &= C_{V,m} + R \\
    C_{p,m} - C_{V,m} &= R \quad \text{(迈耶公式)}
\end{align*}

\section{第三章：输运现象与分子动理学理论的非平衡态理论}
\subsection{黏性现象}
牛顿黏性定律
\begin{align*}
    F &= -\eta \frac{du}{dz} A \\
    J_p &= -\eta \frac{du}{dz}
\end{align*}

重要公式
\begin{align*}
    \frac{dV}{dt} &= \frac{\pi r^4 \Delta p}{8\eta L} \quad \text{(泊肃叶定律)} \\
    f &= 6\pi\eta v R \quad \text{(斯托克斯定律)}
\end{align*}

\subsection{扩散现象}
菲克定律
\begin{align*}
    J_N &= -D \frac{dn}{dz} \\
    \frac{dM}{dt} &= -D \frac{d\rho}{dz} A
\end{align*}

\subsection{热传导现象}
傅里叶定律
\begin{align*}
    \dot{Q} &= -\kappa \frac{dT}{dz} A \\
    J_T &= -\kappa \frac{dT}{dz}
\end{align*}

热欧姆定律
\begin{align*}
    \Delta U_T &= R_T I_T \\
    R_T &= \frac{L}{\kappa A}, \quad \rho_T = \frac{1}{\kappa}
\end{align*}

\subsection{分子平均自由程}
碰撞参数
\begin{align*}
    \sigma &= \pi d^2 \\
    \bar{Z} &= \sqrt{2}n\bar{v}\sigma \\
    \bar{\lambda} &= \frac{\bar{v}}{\bar{Z}} = \frac{1}{\sqrt{2}n\sigma} \\
    \bar{\lambda} &= \frac{kT}{\sqrt{2}\sigma p}
\end{align*}

\subsection{气体输运系数的导出}
微观表达式
\begin{align*}
    \eta &= \frac{1}{3}\rho\bar{v}\bar{\lambda} = \frac{m\bar{v}}{3\sqrt{2}\sigma} \\
    \kappa &= \frac{1}{3}\rho\bar{v}\bar{\lambda} \frac{C_{V,m}}{M_m} \\
    D &= \frac{1}{3}\bar{v}\bar{\lambda}
\end{align*}

\section{第四章：热力学第一定律}
\subsection{可逆与不可逆过程}
准静态过程
中间态均为平衡态的过程

弛豫时间
系统从非平衡态到平衡态所需时间

\subsection{功和热量}
体积功
\begin{align*}
    dW &= -p_edV \\
    dW &= -pdV \quad \text{(准静态过程)} \\
    W &= -\int_{V_1}^{V_2} pdV
\end{align*}

常见过程的功
\begin{align*}
    W &= -vRT\ln\frac{V_2}{V_1} \quad \text{(等温过程)} \\
    W &= -p(V_2 - V_1) \quad \text{(等压过程)} \\
    W &= 0 \quad \text{(等体过程)}
\end{align*}

\subsection{热力学第一定律}
数学表达式
\begin{align*}
    \Delta U &= Q + W \\
    dU &= dQ + dW
\end{align*}

内能定理
绝热过程中 $\Delta U = W_{\text{绝热}}$

\subsection{热容与焓}
热容定义
\begin{align*}
    C_V &= \left(\frac{dQ}{dT}\right)_V, \quad C_{V,m} = \left(\frac{\partial U_m}{\partial T}\right)_V \\
    C_p &= \left(\frac{dQ}{dT}\right)_p, \quad C_{p,m} = \left(\frac{\partial H_m}{\partial T}\right)_p
\end{align*}

焓的定义
\begin{align*}
    H &= U + pV \\
    dH &= dU + pdV + Vdp
\end{align*}

\subsection{第一定律对理想气体的应用}
内能与焓
\begin{align*}
    dU &= vC_{V,m}dT \\
    dH &= vC_{p,m}dT
\end{align*}

过程方程
\begin{align*}
    pV &= \text{常量} \quad \text{(等温)} \\
    pV^\gamma &= \text{常量} \quad \text{(绝热)} \\
    pV^n &= \text{常量} \quad \text{(多方)}
\end{align*}

绝热过程功
\[
W_{\text{绝热}} = \frac{p_2V_2 - p_1V_1}{\gamma - 1} = \frac{vR}{\gamma - 1}(T_2 - T_1)
\]

多方过程
\begin{align*}
    W &= \frac{vR}{n-1}(T_2 - T_1) \\
    C_{n,m} &= C_{V,m} \frac{\gamma - n}{1 - n}
\end{align*}

\subsection{热机与循环}
热机效率
\[
\eta = \frac{W'}{Q_1} = 1 - \frac{|Q_2|}{Q_1}
\]

卡诺热机效率
\[
\eta_{\text{卡}} = 1 - \frac{T_2}{T_1}
\]

内燃机循环
\begin{align*}
    \eta &= 1 - \left(\frac{V_1}{V_2}\right)^{1-\gamma} \quad \text{(奥托循环)} \\
    \eta &= 1 - \frac{\rho^\gamma - 1}{\gamma(\rho - 1)K^{\gamma-1}} \quad \text{(狄塞尔循环)}
\end{align*}

\subsection{焦耳-汤姆孙效应与制冷机}
制冷系数
\[
\eta_{\text{冷}} = \frac{|Q_2|}{W} = \frac{|Q_2|}{|Q_1| - |Q_2|}
\]

卡诺制冷机
\[
\eta_{\text{卡冷}} = \frac{T_2}{T_1 - T_2}
\]

\section{第五章：热力学第二定律与熵}
\subsection{热力学第二定律表述}
开尔文表述：
不能从单一热源吸热完全变功而不产生其他影响

克劳修斯表述：
热量不能自发从低温物体传到高温物体

\subsection{卡诺定理}
定理内容
\begin{itemize}
    \item 可逆热机效率只与热源温度有关
    \item 不可逆热机效率不大于可逆热机
\end{itemize}

卡诺效率
\[
\eta_{\text{可逆}} = 1 - \frac{T_2}{T_1}
\]

\subsection{熵与熵增加原理}
克劳修斯等式
\[
\oint \frac{dQ}{T} = 0 \quad \text{(可逆循环)}
\]

熵的定义
\begin{align*}
    dS &= \frac{dQ_{\text{可逆}}}{T} \\
    \Delta S &= \int \frac{dQ_{\text{可逆}}}{T}
\end{align*}

熵增加原理
\[
(\Delta S)_{\text{绝热}} \geq 0
\]

热力学基本方程
\[
dU = TdS - pdV
\]

理想气体熵变
\begin{align*}
    \Delta S &= \int_{T_1}^{T_2} vC_{V,m}\frac{dT}{T} + vR\ln\frac{V_2}{V_1} \\
    \Delta S &= \int_{T_1}^{T_2} vC_{p,m}\frac{dT}{T} - vR\ln\frac{p_2}{p_1}
\end{align*}

玻尔兹曼熵
\[
S = k\ln W
\]

不可逆过程熵变
\begin{align*}
    \Delta S &> 0 \quad \text{(不可逆过程)} \\
    \Delta S &= vR\ln\frac{V_2}{V_1} \quad \text{(自由膨胀)} \\
    \Delta S &= C_p\ln\frac{(T_1 + T_2)^2}{4T_1T_2} \quad \text{(热传导)}
\end{align*}

\section{常数与单位}
\subsection{基本常数}
\begin{align*}
    R &= 8.31 \, \text{J·mol}^{-1}\text{·K}^{-1} \\
    k &= 1.38 \times 10^{-23} \, \text{J·K}^{-1} \\
    N_A &= 6.02 \times 10^{23} \, \text{mol}^{-1}
\end{align*}

\subsection{标准状况}
\begin{align*}
    T_0 &= 273.15 \, \text{K} \\
    p_0 &= 1 \, \text{atm} = 1.013 \times 10^5 \, \text{Pa} \\
    V_{m,0} &= 22.4 \times 10^{-3} \, \text{m}^3\text{·mol}^{-1}
\end{align*}

\section*{复习建议}
\subsection{重点公式记忆}
建议重点掌握理想气体物态方程、热力学第一定律、各种过程的功和热量计算、热机效率公式、熵变计算公式等。

\subsection{典型题型练习}
多做计算题，特别是涉及多个过程组合的题目，加强对公式应用的理解。

\subsection{概念理解}
重视对基本概念的理解，如平衡态、可逆过程、熵的物理意义等。

\end{document}